\documentclass[11pt,a4paper]{article}

\usepackage[T1]{fontenc}
\usepackage[utf8]{inputenc}
\usepackage[ngerman]{babel}
\usepackage{lmodern}
\usepackage{microtype}
\usepackage[a4paper,margin=2.6cm]{geometry}
\usepackage{amsmath,amssymb,amsthm,mathtools}
\usepackage{booktabs,tabularx,array}
\usepackage{enumitem}
\usepackage{xcolor}
\usepackage{listings}
\usepackage{hyperref}
\usepackage[nameinlink,noabbrev]{cleveref}

\hypersetup{
  pdftitle={Oktonionische Assoziatorgeometrie, triadische Ähnlichkeiten und epistemische Claim Walls},
  pdfauthor={Thomas Hoffbauer},
  colorlinks=true,
  linkcolor=blue!45!black,
  citecolor=blue!45!black,
  urlcolor=blue!45!black
}

\setlist{nosep}
\renewcommand{\arraystretch}{1.15}

\newtheorem{theorem}{Satz}[section]
\newtheorem{proposition}[theorem]{Proposition}
\newtheorem{lemma}[theorem]{Lemma}
\newtheorem{corollary}[theorem]{Korollar}
\theoremstyle{definition}
\newtheorem{definition}[theorem]{Definition}
\newtheorem{remark}[theorem]{Bemerkung}

\newcommand{\Oalg}{\mathbb O}
\newcommand{\ImO}{\operatorname{Im}(\mathbb O)}
\newcommand{\Assoc}{\operatorname{Assoc}}
\newcommand{\MtwoThree}{\mathcal M_{2,3}}
\newcommand{\R}{\mathbb R}
\newcommand{\Nzero}{\mathbb N_0}
\newcommand{\id}{\operatorname{id}}

\lstdefinestyle{logstyle}{
  basicstyle=\ttfamily\scriptsize,
  breaklines=true,
  frame=single,
  columns=fullflexible,
  keepspaces=true,
  backgroundcolor=\color{black!2}
}

\title{\textbf{Oktonionische Assoziatorgeometrie, triadische Ähnlichkeiten\\
und epistemische Claim Walls}\\[0.4em]
\large Ein hochpräzises Audit- und Governance-Modell für \texttt{eabc-renorm}}

\author{
Thomas Hoffbauer\\
Computer Institut Bamberg (CIB)\\
Akademie Olympia\\
\texttt{research@cib-bamberg.de}
}

\date{3. August 2026}

\begin{document}
\maketitle

\begin{abstract}
Wir untersuchen drei miteinander verbundene, aber epistemisch strikt zu trennende Strukturebenen: erstens die exakte Zerlegung des oktonionischen Assoziators auf dem siebendimensionalen Imaginärraum $\ImO$, zweitens die Gram--$G_2$-Identität für seine Norm und drittens triadische Ähnlichkeitstransformationen sowie das explizit definierte dyadisch-triadische Skalierungsmonoid $\MtwoThree$. Für Einheitsvektoren $x,y\in\ImO$ mit $\langle x,y\rangle=-\tfrac12$ tritt in der Vektorraumzerlegung des Assoziators der Summand $\tfrac12 z$ auf. Wir zeigen jedoch, dass dieser Summand keinen longitudinalen Nettoanteil erzeugt, da der vollständige Assoziator orthogonal zu $\operatorname{span}\{x,y,z\}$ ist. Ein Artin-Kontrollfall innerhalb der von $x$ und $y$ erzeugten quaternionischen Unteralgebra zeigt zusätzlich, dass der Winkel $120^\circ$ keine Nichtassoziativität erzwingt. Die Norm des Assoziators wird durch
\[
\|\Assoc(x,y,z)\|^2
=
4\bigl(\det G(x,y,z)-\varphi(x,y,z)^2\bigr)
\]
bestimmt. Für Ähnlichkeiten $\mathcal T_{\lambda,A}=\lambda A$ mit $\lambda=\tfrac13$ werden die Homogenitätsgrade eins, zwei und drei für Längen, quadratische Formen beziehungsweise trilineare Strukturen explizit getrennt. Das Monoid $\MtwoThree=\{2^{-a}3^{-b}\id\}$ wird als exakte mathematische Definition, seine Verwendung als Modellschicht jedoch als revidierbare Modellentscheidung klassifiziert. Ein projektinternes SageMath-Audit auf \texttt{RealField(200)} liefert Residuen in der Größenordnung $10^{-58}$. Der Beitrag formuliert schließlich eine formale Claim Wall, welche die unzulässige Ableitung physikalischer Skalierungsdynamik aus bloßen Zerlegungskoeffizienten ausschließt.
\end{abstract}

\noindent\textbf{Schlagwörter:}
Oktonionen, Assoziator, $G_2$-Geometrie, Malcev-Algebra, Gram-Determinante,
Ähnlichkeitstransformation, glatte Zahlen, mathematische Governance,
numerisches Audit.

\section{Einleitung}

Nichtassoziative Algebren besitzen eine ungewöhnliche Doppelrolle. Einerseits stellen sie präzise und hochstrukturierte mathematische Objekte dar; andererseits laden ihre Symmetrien, Winkelkoeffizienten und Skalengesetze zu physikalischen Deutungen ein, die über die tatsächlich bewiesenen Aussagen hinausgehen können. Im Projekt \texttt{eabc-renorm} wird deshalb eine explizite epistemische Governance verwendet: mathematische Identitäten, numerische Implementierungstests, Modellentscheidungen und physikalische Hypothesen werden unterschiedlichen Aussageklassen zugeordnet.

Der vorliegende Beitrag verfolgt drei Ziele:
\begin{enumerate}
  \item die exakte oktonionische Assoziatorzerlegung unter einer $120^\circ$-Winkelbedingung darzustellen;
  \item die geometrische Norm des Assoziators durch die Gram--$G_2$-Identität zu charakterisieren;
  \item triadische Ähnlichkeiten und das Monoid $\MtwoThree$ einzuführen, ohne einen Zerlegungskoeffizienten fälschlich als dynamischen Skalierungsoperator zu interpretieren.
\end{enumerate}

Die numerischen Resultate werden als projektinterne, reproduzierbare Auditdaten berichtet. Sie ersetzen keinen mathematischen Beweis, sondern prüfen Implementierung, Vorzeichenkonventionen und numerische Konsistenz.

\section{Algebraische Grundlagen}

\subsection{Oktonionen und Kreuzprodukt}

Wir schreiben
\[
\Oalg=\R\oplus\ImO,
\qquad
\ImO\cong\R^7.
\]
Für $x,y\in\ImO$ gilt
\begin{equation}
xy=-\langle x,y\rangle+x\times y,
\label{eq:imag-product}
\end{equation}
wobei $x\times y\in\ImO$ das siebendimensionale Kreuzprodukt bezeichnet. Es ist alternierend und erfüllt
\[
\|x\times y\|^2
=
\|x\|^2\|y\|^2-\langle x,y\rangle^2.
\]
Die $G_2$-Dreiform ist
\begin{equation}
\varphi(x,y,z)
=
\langle x\times y,z\rangle.
\label{eq:g2-form}
\end{equation}

\begin{definition}[Assoziator]
Für $x,y,z\in\Oalg$ sei
\begin{equation}
\Assoc(x,y,z)
\coloneqq
(xy)z-x(yz).
\end{equation}
\end{definition}

Da $\Oalg$ alternativ ist, ist der Assoziator alternierend. Für rein imaginäre Argumente ist er wiederum rein imaginär und orthogonal zu jedem seiner drei Argumente.

\subsection{Exakte Vektorraumzerlegung}

\begin{theorem}[Assoziatorzerlegung]
\label{thm:assoc-decomp}
Für $x,y,z\in\ImO$ gilt
\begin{equation}
\Assoc(x,y,z)
=
-\langle x,y\rangle z
+\langle y,z\rangle x
+(x\times y)\times z
-x\times(y\times z).
\label{eq:assoc-decomp}
\end{equation}
\end{theorem}

\begin{proof}
Aus \cref{eq:imag-product} folgt
\begin{align}
(xy)z
&=
-\langle x\times y,z\rangle
-\langle x,y\rangle z
+(x\times y)\times z,\\
x(yz)
&=
-\langle x,y\times z\rangle
-\langle y,z\rangle x
+x\times(y\times z).
\end{align}
Wegen
\[
\langle x\times y,z\rangle
=
\langle x,y\times z\rangle
=
\varphi(x,y,z)
\]
heben sich die skalaren Dreifachprodukte in der Differenz auf.
\end{proof}

\subsection{Der $120^\circ$-Koeffizient und seine Kompensation}

\begin{proposition}[Winkelterm ohne longitudinalen Nettoanteil]
\label{prop:120}
Seien $x,y\in\ImO$ Einheitsvektoren mit
\[
\langle x,y\rangle=\cos(120^\circ)=-\frac12.
\]
Dann lautet der erste Summand in \cref{eq:assoc-decomp}
\[
-\langle x,y\rangle z=\frac12z.
\]
Gleichzeitig gilt für alle $z\in\ImO$
\[
\langle\Assoc(x,y,z),z\rangle=0,
\]
also
\begin{equation}
\left\langle
\langle y,z\rangle x
+(x\times y)\times z
-x\times(y\times z),
z
\right\rangle
=
-\frac12\|z\|^2.
\label{eq:compensation}
\end{equation}
\end{proposition}

\begin{remark}
Der Term $\tfrac12z$ ist ein winkelabhängiger Bestandteil einer Zerlegung. Er ist weder ein isolierter Anteil des Gesamtassoziators noch eine lineare Dynamik $z\mapsto z/2$.
\end{remark}

\subsection{Artin-Kontrollfall}

\begin{proposition}[Assoziative Zweierunteralgebra]
\label{prop:artin}
Sei
\[
\mathbb A_{x,y}
=
\operatorname{alg}_{\R}\langle x,y\rangle
\subset\Oalg.
\]
Nach dem Satz von Artin ist $\mathbb A_{x,y}$ assoziativ. Sind $x$ und $y$ rein imaginär, normiert und linear unabhängig, so gilt
\[
\mathbb A_{x,y}
=
\operatorname{span}_{\R}\{1,x,y,x\times y\}
\cong_{\R\text{-alg}}\mathbb H.
\]
Für jedes
\[
z\in\mathbb A_{x,y}\cap\ImO
\]
folgt
\[
\Assoc(x,y,z)=0.
\]
\end{proposition}

Damit verschwindet der Assoziator beispielsweise für
\[
z_{\mathbb A}
=
\frac{x\times y}{\|x\times y\|},
\]
obwohl der Summand $\tfrac12z_{\mathbb A}$ in der Zerlegung weiterhin vorkommt. Der Winkelkoeffizient ist folglich kein Maß der Nichtassoziativität.

\section{Gram--$G_2$-Geometrie}

Für $v_1=x$, $v_2=y$, $v_3=z$ sei
\[
G(x,y,z)
=
\bigl(\langle v_i,v_j\rangle\bigr)_{i,j=1}^{3}
\]
die Gram-Matrix.

\begin{theorem}[Gram--$G_2$-Identität]
\label{thm:gram-g2}
Für alle $x,y,z\in\ImO$ gilt
\begin{equation}
\|\Assoc(x,y,z)\|^2
=
4\left(
\det G(x,y,z)-\varphi(x,y,z)^2
\right).
\label{eq:gram-g2}
\end{equation}
Insbesondere
\[
0\leq \varphi(x,y,z)^2\leq \det G(x,y,z).
\]
\end{theorem}

Die Differenz
\[
\det G-\varphi^2
\]
misst den Anteil des euklidischen Dreiervolumens, der nicht durch die $G_2$-kalibrierte Dreiform erfasst wird. \Cref{eq:gram-g2} bestimmt die Norm, nicht die vollständige Richtung des Assoziators.

\subsection{Jacobiator}

Definiere
\[
J_\times(x,y,z)
=
(x\times y)\times z
+(y\times z)\times x
+(z\times x)\times y.
\]

\begin{theorem}[Jacobiator-Assoziator-Relation]
Für die Konventionen
\[
\Assoc(x,y,z)=(xy)z-x(yz),
\qquad
x\times y=\frac12(xy-yx)
\]
gilt
\begin{equation}
J_\times(x,y,z)
=
\frac32\Assoc(x,y,z).
\label{eq:jac-assoc}
\end{equation}
\end{theorem}

\begin{proof}
Für den Kommutator $[x,y]_{\mathrm c}=xy-yx$ gilt in einer alternativen Algebra
\[
[[x,y]_{\mathrm c},z]_{\mathrm c}
+
[[y,z]_{\mathrm c},x]_{\mathrm c}
+
[[z,x]_{\mathrm c},y]_{\mathrm c}
=
6\Assoc(x,y,z).
\]
Jeder doppelt geschachtelte Kreuzproduktterm ist ein Viertel des entsprechenden Kommutatorterms. Daraus folgt \cref{eq:jac-assoc}.
\end{proof}

\section{Triadische Ähnlichkeitstransformationen}

\subsection{Dreidimensionaler Trägerraum und globale Wirkung}

Sei $V\subset\ImO$ ein dreidimensionaler euklidischer Unterraum und $\Delta\subset V$ ein reguläres Tetraeder. Weiter sei
\[
A\in O(\ImO),
\qquad
A(V)=V,
\]
und $A_V=A|_V$.

\begin{definition}[Triadische Ähnlichkeit]
Für
\[
\lambda=\frac13
\]
definieren wir
\[
\mathcal T_{\lambda,A}
=
\lambda A.
\]
\end{definition}

Dann gilt
\[
\mathcal T_{\lambda,A}^{\,k}
=
3^{-k}A^k.
\]
Kantenlängen skalieren mit $3^{-k}$, dreidimensionale Volumina mit $27^{-k}$.

Für
\[
N(x)=\|x\|^2,
\qquad
Q(x)=\frac12\|x\|^2
\]
gilt
\begin{equation}
N(\mathcal T_{\lambda,A}x)
=
\frac19N(x),
\qquad
Q(\mathcal T_{\lambda,A}x)
=
\frac19Q(x).
\label{eq:quadratic-scale}
\end{equation}

\subsection{Rotation und zentrale Inversion}

Eine Halbdrehung um die Achse $\R n\subset V$ ist
\[
R_{\pi,n}x
=
2\langle n,x\rangle n-x.
\]
Sie besitzt in drei Dimensionen die Eigenwerte
\[
1,-1,-1
\]
und Determinante $+1$. Die zentrale Inversion $-I_V$ besitzt dagegen die Eigenwerte
\[
-1,-1,-1
\]
und Determinante $-1$. Daher
\[
R_{\pi,n}\neq -I_V.
\]

Nur für die globale zentrale Inversion $A=-I_7$ gilt
\[
\mathcal T_{\lambda,-I_7}^{\,k}x
=
\left(-\frac13\right)^kx.
\]

\subsection{Homogenitätsgrade}

Für jede trilineare Abbildung $F$ gilt
\[
F(\lambda x,\lambda y,\lambda z)
=
\lambda^3F(x,y,z).
\]
Somit skalieren
\[
\varphi,\qquad \Assoc,\qquad J_\times
\]
bei gemeinsamer Skalierung ihrer drei Argumente mit $\lambda^3$.

Für die $G_2$-Dreiform gilt allgemeiner
\begin{equation}
\varphi(\lambda Ax,\lambda Ay,\lambda Az)
=
\lambda^3(A^*\varphi)(x,y,z),
\end{equation}
wobei
\[
(A^*\varphi)(x,y,z)
=
\varphi(Ax,Ay,Az).
\]

Ist $A\in G_2$, dann
\[
A^*\varphi=\varphi,
\qquad
A(x\times y)=Ax\times Ay,
\]
und folglich
\[
J_\times(\lambda Ax,\lambda Ay,\lambda Az)
=
\lambda^3A\,J_\times(x,y,z).
\]

Für $A=-I_7$ gilt dagegen
\[
A^*\varphi=-\varphi.
\]

\section{Das dyadisch-triadische Skalierungsmonoid}

\begin{definition}
Auf einem reellen Vektorraum seien
\[
D(x)=\frac12x,
\qquad
T(x)=\frac13x.
\]
Das von $D$ und $T$ erzeugte kommutative Monoid ist
\begin{equation}
\MtwoThree
=
\{D^aT^b\mid a,b\in\Nzero\}
=
\{2^{-a}3^{-b}\id\mid a,b\in\Nzero\}.
\label{eq:monoid}
\end{equation}
\end{definition}

Es gilt
\[
D^aT^b=T^bD^a
\]
und für jede quadratische Form
\begin{equation}
Q(D^aT^bx)
=
4^{-a}9^{-b}Q(x).
\label{eq:monoid-q}
\end{equation}

Die mathematische Definition von $\MtwoThree$ besitzt den Rang \texttt{[A-def]}. Die Auswahl gerade dieses Monoids als Modellschicht von \texttt{eabc-renorm} ist dagegen eine revidierbare Modellentscheidung \texttt{[B-model]}.

\section{Numerisches Hochpräzisionsaudit}

\subsection{Methodik}

Nach projektinternem Auditprotokoll wurden zwei SageMath-Suiten eingesetzt:
\begin{enumerate}
  \item eine konventionskalibrierte, implementierungsentkoppelte Fano-Rekonstruktion der Assoziatorzerlegung;
  \item eine getrennte Prüfung dreidimensionaler Transformationsidentitäten und siebendimensionaler Homogenitätsgesetze.
\end{enumerate}

Die linke Seite der Assoziatoridentität wurde über die native Oktonionenmultiplikation berechnet. Die rechte Seite wurde über einen explizit implementierten Fano-Tensor rekonstruiert. Die beiden Auswertungspfade teilen dieselbe Oktonionenstruktur und stellen daher keinen zweiten mathematischen Beweis dar; sie sind jedoch als Implementierungspfade entkoppelt.

Als numerische Umgebung wurden nach dem vorgelegten Protokoll verwendet:
\[
\texttt{SageMath 10.5},
\qquad
\texttt{RealField(200)},
\qquad
\texttt{samples}=100,
\qquad
\varepsilon_{\mathrm{tol}}=10^{-45}.
\]

\subsection{Berichtete Resultate}

\begin{table}[ht]
\centering
\small
\caption{Projektintern berichtete Auditwerte. Die Werte sind als Provenienzdaten des Repository-Audits zu verstehen.}
\label{tab:audit}
\begin{tabularx}{\textwidth}{@{}Xrr@{}}
\toprule
\textbf{Prüfdimension} & \textbf{Residuum/Wert} & \textbf{Status}\\
\midrule
Fano-Basiskalibrierung & $0$ & \texttt{PASS}\\
Unprojizierte Reinheit des 8D-Assoziators & $0$ & \texttt{PASS}\\
$120^\circ$-Winkelbedingung & $0$ & \texttt{PASS}\\
Kreuzproduktnorm & $1{,}9721\cdot10^{-59}$ & \texttt{PASS}\\
Fano-Assoziatorrekonstruktion & $3{,}9443\cdot10^{-59}$ & \texttt{PASS}\\
Volle Orthogonalität & $1{,}9721\cdot10^{-59}$ & \texttt{PASS}\\
Kompensationsidentität & $1{,}9721\cdot10^{-59}$ & \texttt{PASS}\\
Alternierung & $0$ & \texttt{PASS}\\
Antisymmetrie & $0$ & \texttt{PASS}\\
Artin-Nullkontrolle & $0$ & \texttt{PASS}\\
Gram--$G_2$-Identität & $1{,}5777\cdot10^{-58}$ & \texttt{PASS}\\
Minimal beobachteter Gap $\det G-\varphi^2$ & $1{,}4712\cdot10^{-3}$ & \texttt{PASS}\\
$\det R_{\pi,n}=+1$ in $3$D & $0$ & \texttt{PASS}\\
$\det(-I_3)=-1$ & $0$ & \texttt{PASS}\\
Achsenunterscheidung & $0$ & \texttt{PASS}\\
Kubische Jacobiatorskalierung & $1{,}9721\cdot10^{-59}$ & \texttt{PASS}\\
Monoidkommutativität & $0$ & \texttt{PASS}\\
Quadratische Monoidwirkung & $0$ & \texttt{PASS}\\
\bottomrule
\end{tabularx}
\end{table}

Als Provenienz wurden angegeben:
\begin{itemize}
  \item Assoziator-Audit: Commit \texttt{7ff7c62d8e41a8b9c204918f3a5619d08e7b4c91};
  \item triadisches Modul: Commit \texttt{8dbfd66};
  \item Protokolldatei: \texttt{artifacts/audit\_120deg\_assoc\_RF200.txt}.
\end{itemize}

Die genannten Commit- und Mergeangaben werden hier als vom Projekt bereitgestellte Provenienzinformation wiedergegeben.

\section{Epistemische Governance}

\begin{table}[ht]
\centering
\small
\caption{Epistemische Klassifikation der behandelten Aussagen.}
\begin{tabularx}{\textwidth}{@{}p{0.20\textwidth}p{0.17\textwidth}X@{}}
\toprule
\textbf{Kategorie} & \textbf{Rang} & \textbf{Geltungsbereich}\\
\midrule
Algebraische Identitäten & \texttt{[A-kernel]} &
Assoziatorzerlegung, Artin-Kontrolle, Gram--$G_2$-Identität,
Jacobiatorrelation und Homogenitätsgrade.\\
Monoidstruktur & \texttt{[A-def]} &
Exakte Definition von $\MtwoThree$ und seine Wirkung auf Vektoren
und quadratische Formen.\\
Modellwahl & \texttt{[B-model]} &
Auswahl von $\MtwoThree$ als Skalierungsarchitektur von
\texttt{eabc-renorm}.\\
Numerisches Audit & \texttt{[B-audit]} &
Hochpräzise Regressionstests der Implementierung und Konventionen.\\
Claim Wall & \texttt{NON-CLAIM} &
Keine Herleitung von $D=\tfrac12I$ aus dem Zerlegungsterm
$\tfrac12z$.\\
Physikalische Dynamik & \texttt{[C]/NON-CLAIM} &
Raumzeitquantisierung, Eichfelder oder Renormierungsflüsse erfordern
zusätzliche Dynamik und Evidenz.\\
\bottomrule
\end{tabularx}
\end{table}

\subsection{Formale Claim Wall}

Die zulässige und die unzulässige Schlusskette lassen sich als
\[
120^\circ
\longrightarrow
-\langle x,y\rangle z=\frac12z
\not\Longrightarrow
D=\frac12I
\]
und
\[
D=\frac12I,\quad T=\frac13I
\longrightarrow
\MtwoThree
\not\Longrightarrow
\text{Raumzeitquantisierung}
\]
darstellen.

\begin{quote}
\textbf{Claim Wall.}
Der Faktor $\tfrac12$ in der $120^\circ$-Assoziatorzerlegung ist ein
winkelabhängiger Koeffizient einer algebraischen Identität. Er
definiert keine dyadische Homothetie. Das Monoid $\MtwoThree$ entsteht
erst durch die explizite Definition der Operatoren $D$ und $T$. Seine
Verwendung als Modellschicht begründet weder einen physikalischen
Renormierungsfluss noch eine Raumzeitquantisierung.
\end{quote}

\section{Offene Forschungsfragen}

Nach Schließung der internen Konsistenzfragen verbleibt die eigentliche Modellfrage:
\begin{equation}
\boxed{
\text{Warum sollte }\MtwoThree
\text{ eine ausgezeichnete Skalierungsstruktur für EABC sein?}
}
\end{equation}

Mögliche Untersuchungsrichtungen sind:
\begin{enumerate}
  \item spektrale Eigenschaften diskreter Gitterprojektionen und die Frage, ob $2$- und $3$-glatte Skalen ausgezeichnete Eigen- oder Selbstähnlichkeitsstrukturen erzeugen;
  \item rationale und $3$-adische Modellräume, in denen
  \[
  \left\|\frac13v\right\|_3=3\|v\|_3
  \]
  der archimedischen Kontraktion gegenübersteht;
  \item algebraische Wirkungen von Moufang-Schleifen auf explizit definierten Selbstähnlichkeitsoperatoren;
  \item Vergleich des Monoids $\MtwoThree$ mit alternativen Skalenmonoiden und Nullmodellen;
  \item Entwicklung einer eigenständigen Dynamik, bevor physikalische Aussagen über Renormierung oder Raumzeit getroffen werden.
\end{enumerate}

\section{Schlussfolgerung}

Die mathematische Analyse trennt einen Winkelkoeffizienten in einer oktonionischen Zerlegung von einer eigenständigen Skalierungsoperation. Der Term $\tfrac12z$ besitzt keine longitudinale Nettowirkung im Assoziator und erzwingt keine Nichtassoziativität. Die Gram--$G_2$-Identität liefert ein geometrisch präzises Normmaß der nichtassoziativen Abweichung. Triadische Ähnlichkeiten skalieren Längen, quadratische Formen und trilineare Strukturen entsprechend ihren Homogenitätsgraden. Das Monoid $\MtwoThree$ ist mathematisch wohldefiniert, seine Modellrolle jedoch nicht aus der Oktonionenalgebra hergeleitet.

Die zentrale methodische Folgerung lautet:
\[
\boxed{
\text{Strukturidentität}
\neq
\text{Modelloperator}
\neq
\text{physikalische Dynamik}.
}
\]

\appendix

\section{Abfrage und Ereignisprotokoll}

Die projektinterne Abfrage für Claim-Wall-Verifikationsereignisse lautet:
\begin{lstlisting}[style=logstyle]
type:SecurityEvent
| where Category == "AlgebraicGovernance"
  AND Action == "ClaimWallVerification"
\end{lstlisting}

Ein gekürzter CEF-Ereignisstrom wurde in folgender Form berichtet:
\begin{lstlisting}[style=logstyle]
CEF:0|CIB|eabc-renorm|10.5|GOV-101|Octonionic Assoc Vector Decomp|1|
msg=Exact vector space decomposition confirmed for 343 basis triples.

CEF:0|CIB|eabc-renorm|10.5|GOV-103|Longitudinal Orthogonality Guard|1|
msg=Orthogonality <Assoc(x,y,z),z>=0 verified.

CEF:0|CIB|eabc-renorm|10.5|GOV-105|Gram-G2 Volume Identity|1|
msg=||Assoc||^2 = 4*(det(G)-phi^2) confirmed.

CEF:0|CIB|eabc-renorm|10.5|GOV-201|Triadic Homothetia Degree-3 Check|1|
msg=phi, Assoc and J_cross scale cubically with lambda^3=1/27.

CEF:0|CIB|eabc-renorm|10.5|GOV-202|3D Transformation Distinction|1|
msg=det(R_pi_3D)=+1 and det(-I_3D)=-1 verified.

CEF:0|CIB|eabc-renorm|10.5|GOV-203|Monoid M_2,3 Commutativity|1|
msg=D^a T^b == T^b D^a confirmed.
\end{lstlisting}

\section*{Hinweis zur Reproduzierbarkeit}

Die mathematischen Aussagen des Artikels sind von den projektintern
berichteten Audit- und Repository-Metadaten zu unterscheiden. Die
numerischen Werte, Commit-Kennungen und Mergeangaben werden aus dem
bereitgestellten Laborbericht übernommen. Für eine arXiv-Einreichung
sollten Quellskripte, vollständige Commit-Hashes, Maschinenprotokolle
und eine eindeutige Softwareumgebung dauerhaft öffentlich archiviert
werden.

\begin{thebibliography}{9}

\bibitem{Baez2002}
J.~C. Baez,
\newblock The Octonions,
\newblock \emph{Bulletin of the American Mathematical Society}
39 (2002), 145--205.

\bibitem{Harvey1990}
F.~R. Harvey,
\newblock \emph{Spinors and Calibrations},
\newblock Academic Press, Boston, 1990.

\bibitem{Schafer1966}
R.~D. Schafer,
\newblock \emph{An Introduction to Nonassociative Algebras},
\newblock Academic Press, New York, 1966.

\bibitem{SpringerVeldkamp2000}
T.~A. Springer und F.~D. Veldkamp,
\newblock \emph{Octonions, Jordan Algebras and Exceptional Groups},
\newblock Springer, Berlin, 2000.

\end{thebibliography}

\end{document}
